\documentclass[11pt,a4paper]{article}

% Packages
\usepackage[utf8]{inputenc}
\usepackage[T1]{fontenc}
\usepackage{amsmath,amssymb,amsfonts}
\usepackage{graphicx}
\usepackage{hyperref}
\usepackage{natbib}
\usepackage{booktabs}
\usepackage{algorithm}
\usepackage{algorithmic}
\usepackage[margin=1in]{geometry}

% Metadata
\title{Your Paper Title Here}
\author{
  Author One\thanks{Affiliation One} \\
  \texttt{author1@email.com}
  \and
  Author Two\thanks{Affiliation Two} \\
  \texttt{author2@email.com}
}
\date{\today}

\begin{document}

\maketitle

\begin{abstract}
  A concise summary of the paper: the problem addressed, the proposed approach,
  key results, and implications for the Gaussia framework.
\end{abstract}

\section{Introduction}

Motivate the problem. Why does AI evaluation need this metric or approach?
What gap does it fill in the current Gaussia framework?

\section{Related Work}

Discuss existing approaches and how your proposal differs or improves upon them.

\section{Methodology}

Describe your approach in detail. Include:
\begin{itemize}
  \item Formal definitions and notation
  \item Algorithms (use the \texttt{algorithm} environment)
  \item Mathematical formulations
\end{itemize}

\subsection{Formal Definition}

% Example: define your metric formally
% \begin{equation}
%   M(x) = \frac{1}{N} \sum_{i=1}^{N} f(x_i)
% \end{equation}

\subsection{Algorithm}

% Example algorithm
% \begin{algorithm}
%   \caption{Your Algorithm}
%   \begin{algorithmic}[1]
%     \REQUIRE Input data $X$
%     \ENSURE Metric score $M$
%     \STATE Initialize $M \leftarrow 0$
%     \FOR{each $x_i \in X$}
%       \STATE $M \leftarrow M + f(x_i)$
%     \ENDFOR
%     \RETURN $M / |X|$
%   \end{algorithmic}
% \end{algorithm}

\section{Implementation Considerations}

Describe how this should be implemented in the Gaussia framework:
\begin{itemize}
  \item Which base class to extend (\texttt{Gaussia}, \texttt{Guardian}, etc.)
  \item Required dependencies
  \item Expected schema for results
  \item Performance considerations
  \item Statistical mode compatibility (Frequentist/Bayesian)
\end{itemize}

\section{Evaluation}

If applicable, present experimental results:
\begin{itemize}
  \item Datasets used
  \item Baselines compared against
  \item Results and analysis
\end{itemize}

% Example table
% \begin{table}[h]
%   \centering
%   \caption{Experimental Results}
%   \begin{tabular}{lcc}
%     \toprule
%     Method & Score & Runtime \\
%     \midrule
%     Baseline & 0.72 & 1.2s \\
%     Ours & \textbf{0.89} & 0.8s \\
%     \bottomrule
%   \end{tabular}
% \end{table}

\section{Conclusion}

Summarize contributions and outline future work.

\bibliographystyle{plainnat}
\bibliography{references}

\end{document}
